%%%%%%
%
% $Autor: Wings $
% $Datum: 2020-01-18 11:15:45Z $
% $Pfad: githubtemplate/Template/Presentations/Template/slides/rename.tex $
% $Version: 4620 $
%
%
% !TeX encoding = utf8
% !TeX root = Rename
%
%%%%%%

	
	\Mysection{Introduction}
	\STANDARD{Problem Definition}{
		Ground robots require accurate localization to navigate and map environments autonomously.  
		This work evaluates the suitability of:
		\begin{itemize}
			\item \textbf{Visual SLAM algorithms} (e.g., ORB-SLAM2, ORB-SLAM3), and
			\item \textbf{IRSNet-based CNN approaches} for acoustic localization.
		\end{itemize}
		The platform under consideration uses:
		\begin{itemize}
			\item \textbf{NVIDIA Jetson Nano Developer Kit} — embedded GPU for edge AI tasks.
			\item \textbf{Raspberry Pi Camera Module V2 (Sony IMX219, 8 MP)} — RGB input for visual odometry.
		\end{itemize}
	}
	
	\Mysection{Hardware Setup}
	\STANDARD{System Overview}{
		\begin{itemize}
			\item Jetson Nano provides GPU acceleration for real-time SLAM.
			\item Raspberry Pi Camera V2 provides monocular RGB images.
			\item Lightweight robotics middleware (e.g., ROS Noetic) for data flow.
		\end{itemize}
		\vspace{2mm}
	\begin{center}
		\begin{tikzpicture}[
			node distance=6mm,
			every node/.style={font=\footnotesize},
			scale=0.7, transform shape
			]
			% Node and arrow styles
			\tikzstyle{block}=[rectangle, draw, rounded corners,
			minimum width=3.2cm, minimum height=0.8cm, text centered]
			\tikzstyle{arrow}=[->, thick, >=Stealth]
			
			% Nodes (top to bottom)
			\node[block] (jet) {Jetson Nano Platform};
			\node[block, below=of jet] (cam) {Raspberry Pi Camera V2};
			\node[block, below=of cam] (feat) {Feature Extraction (ORB)};
			\node[block, below=of feat] (slam) {SLAM Mapping \& Pose Estimation};
			\node[block, below=of slam] (loc) {Ground Localization Output};
			
			% Arrows
			\draw[arrow] (jet) -- (cam);
			\draw[arrow] (cam) -- (feat);
			\draw[arrow] (feat) -- (slam);
			\draw[arrow] (slam) -- (loc);
		\end{tikzpicture}
	\end{center}
	
	}
	
	\Mysection{SLAM-Based Localization}
	\STANDARD{Feature-Based Visual SLAM}{
		\textbf{ORB-SLAM2} and \textbf{ORB-SLAM3} are state-of-the-art real-time visual SLAM systems.
		\begin{itemize}
			\item Detect and track ORB features across frames.
			\item Estimate camera pose via bundle adjustment.
			\item Map 3D landmarks for loop closure and re-localization.
		\end{itemize}
		\vspace{2mm}
		Performance on Jetson Nano has been validated by Krško et al.\ (2025)\cite{Krsko:2025},  
		demonstrating reliable real-time localization with the Raspberry Pi Camera V2.
		\begin{itemize}
			\item \cite{MurArtal:2017} — ORB-SLAM2 core algorithm.  
			\item \cite{Campos:2021} — ORB-SLAM3 visual-inertial extension.
		\end{itemize}
	}
	
	\STANDARD{Active SLAM Approaches}{
		Active SLAM algorithms plan robot motion to reduce map uncertainty while localizing.  
		Ahmed et al.\ (2023)\cite{Ahmed:2023} review key advances integrating exploration and mapping.  
		However, computational complexity makes them less suitable for real-time embedded deployment on Jetson Nano without simplification.
	}
	
	\Mysection{IRSNet-Based CNN Localization}
	\STANDARD{IRSNet Overview}{
		\textbf{IRSNet}\cite{Sun:2024} is a convolutional neural network for \textit{acoustic} source localization.  
		It estimates sound source positions using spectrogram features and improved residual structures.
		\vspace{2mm}
		\textbf{Limitation for Ground Localization:}
		\begin{itemize}
			\item Requires microphone arrays, not camera data.
			\item Lacks geometric consistency for visual odometry.
			\item Heavy inference cost on Jetson Nano compared to ORB-SLAM.
		\end{itemize}
		Thus, IRSNet is not optimal for vision-based robot localization.
	}
	
	\Mysection{Comparative Evaluation}
	\STANDARD{Algorithm Comparison}{
		\begin{table}[h!]
			\centering
		\resizebox{0.9\textwidth}{!}{
			\begin{tabular}{|l|c|c|c|}
				\hline
				\textbf{Method} & \textbf{Sensor} & \textbf{Accuracy} & \textbf{Jetson Nano Suitability} \\ \hline
				ORB-SLAM2 \cite{MurArtal:2017} & Camera (RGB) & High (cm-level) & Excellent (20–30 fps) \\ \hline
				ORB-SLAM3 \cite{Campos:2021} & Camera + IMU & Very High / Robust & Good (15 fps) \\ \hline
				Active SLAM \cite{Ahmed:2023} & Camera + Depth & High / Exploratory & Moderate (CPU intensive) \\ \hline
				IRSNet \cite{Sun:2024} & Microphone Array & Acoustic only & Not applicable \\ \hline
			\end{tabular}
		}
			\caption{Comparison of Localization Methods for Jetson Nano + Raspberry Pi Camera V2.}
		\end{table}
	
	}
	
	\Mysection{Flowchart}
	\STANDARD{Ground Localization Pipeline}{
		\begin{center}
			\begin{tikzpicture}[
				node distance=8mm,
				every node/.style={font=\footnotesize},
				scale=0.8, transform shape
				]
				\tikzstyle{proc}=[rectangle, draw, rounded corners,
				minimum width=3.4cm, minimum height=0.8cm, text centered, fill=gray!10]
				\tikzstyle{arr}=[->, thick, >=Stealth]
				
				\node[proc] (cap) {Image Capture (RPi Camera V2)};
				\node[proc, below=of cap] (feat) {Feature Extraction (ORB)};
				\node[proc, below=of feat] (pose) {Pose Estimation + Mapping (SLAM)};
				\node[proc, below=of pose] (map) {Localization / Control Loop Output};
				
				\draw[arr] (cap) -- (feat);
				\draw[arr] (feat) -- (pose);
				\draw[arr] (pose) -- (map);
			\end{tikzpicture}
		\end{center}
		
		
	}
	
	\Mysection{Conclusion}
	\STANDARD{Summary and Recommendation}{
		\begin{itemize}
			\item \textbf{IRSNet} excels in acoustic tasks, but not visual localization.
			\item \textbf{ORB-SLAM2 / ORB-SLAM3} are proven, efficient, and open-source for vision-based ground robots.
			\item Optimizations for Jetson Nano: reduce frame resolution, enable loop closure at intervals, use GPU-accelerated features.
			\item Future work: integrate lightweight CNN depth estimation with ORB-SLAM3 for enhanced floor-plane awareness.
		\end{itemize}
	}
	
	\Mysection{References}
	\STANDARD{Key References}{
		\begin{itemize}
			\item Mur-Artal R. and Tardós J.D., “ORB-SLAM2,” IEEE TRO (2017)\cite{MurArtal:2017}
			\item Campos C. et al., “ORB-SLAM3,” IEEE TRO (2021)\cite{Campos:2021}
			\item Ahmed M.F. et al., “Active SLAM: A Review,” Sensors (2023)\cite{Ahmed:2023}
			\item Sun D. et al., “IRSNet,” Remote Sensing (2024)\cite{Sun:2024}
			\item Krško J. et al., “Evaluation of ORB-SLAM2 on Jetson Nano,” Sensors (2025)\cite{Krsko:2025}
		\end{itemize}
	}


